\documentclass[11pt,a4paper,spanish]{book}
\usepackage{estilo_unir-1}
\usepackage{amsmath}
\hypersetup{
	colorlinks   = true, %Colours links instead of ugly boxes
	urlcolor     = blue, %Colour for external hyperlinks
	linkcolor    = black, %Colour of internal links
	citecolor   = red %Colour of citations
}
\usepackage[nottoc]{tocbibind}

%---------------------------
%título de la asignatura, trabajo y autor
%---------------------------
\subject{Mecánica Cuántica}
\title{Dinámica y constantes de movimiento en sistemas hermíticos}
\author{Enrique Prada Vázquez}
\profesor{Rodrigo Gil-Merino y Rubio}
\date{13-04-2026}

%---------------------------
%marges
%---------------------------
%\usepackage[margin=1.9cm]{geometry}
%---------------------------
%---------------------------
%---------------------------
%---------------------------
\begin{document}
\renewcommand{\listfigurename}{Índice de figuras}
\renewcommand{\listtablename}{Índice de tablas}
\renewcommand{\contentsname}{Índice de contenidos}
\renewcommand{\figurename}{Figura}
\renewcommand{\tablename}{Tabla} 

\maketitle
\frontmatter
\tableofcontents

\chapter{Resumen}
Este Trabajo presenta un análisis de un sistema cuántico unidimensional en un espacio de Hilbert de dimensión $d=4$. El propósito principal es caracterizar la dinámica y las propiedades estadísticas del sistema a partir de un Hamiltoniano específico. La metodología se fundamenta en la resolución de la ecuación de autovalores para determinar el espectro energético y la construcción de una base ortonormal de autoestados. Se aborda el fenómeno de la degeneración en el nivel de energía $E=2$ y se realiza la descomposición espectral del operador mediante proyectores. 

Posteriormente, se estudia la evolución temporal de un estado inicial en superposición, aplicando el principio de superposición y la separación de variables según la ecuación de Schrödinger. Los resultados demuestran que, bajo la interpretación probabilística de Max Born, las probabilidades de ocupación de ciertos autoestados permanecen constantes en el tiempo. Finalmente, se investiga la relación entre constantes de movimiento y simetrías, identificando observables compatibles que conmutan con el Hamiltoniano.

\vspace{0.5cm}
{\bf Palabras clave:} mecánica cuántica, espacio de hilbert, hamiltoniano, evolución temporal, regla de born, constantes de movimiento

\chapter{Abstract}
This Report provides an analysis of a one-dimensional quantum system within a four-dimensional Hilbert space. The primary objective is to characterize the dynamics and statistical properties of the system based on a specific Hamiltonian. The methodology is centered on solving the eigenvalue equation to determine the energy spectrum and constructing an orthonormal basis of eigenstates. The work addresses the phenomenon of degeneracy at the $E=2$ energy level and performs the spectral decomposition of the operator using projectors.

Furthermore, the temporal evolution of an initial superposition state is studied by applying the superposition principle and the separation of variables according to the Schrödinger equation. The results demonstrate that, under Max Born's probabilistic interpretation, the occupation probabilities of certain eigenstates remain constant over time. Finally, the relationship between constants of motion and symmetries is investigated, identifying compatible observables that commute with the Hamiltonian.

\vspace{0.5cm}
{\bf Keywords:} quantum mechanics, hilbert space, hamiltonian, time evolution, born rule, constants of motion

\mainmatter


\chapter{Introducción}

La mecánica cuántica representa uno de los pilares fundamentales de la física moderna, proporcionando el marco teórico necesario para comprender el comportamiento de la materia a escalas microscópicas \citep{planck1900}. A diferencia de la mecánica clásica, donde el estado de un sistema se define por variables deterministas, el formalismo cuántico se sustenta en una estructura algebraica de espacios vectoriales complejos y operadores lineales \citep{dirac1930}. El presente Trabajo se enmarca en este contexto, analizando la dinámica de un sistema mediante herramientas fundamentales como la descomposición espectral y la evolución temporal.

\section{Motivación y justificación}

La motivación de este Trabajo radica en la importancia de comprender cómo la energía, representada por el operador Hamiltoniano, condiciona la evolución y la medición de los sistemas físicos. El Hamiltoniano no solo determina los niveles de energía permitidos a través de sus autovalores, sino que también dicta la estabilidad del sistema mediante sus autoestados \citep{schrodinger1926}. 

El estudio de sistemas en espacios de Hilbert de dimensión finita es esencial en campos de vanguardia como la computación cuántica y la información cuántica \citep{feynman1982}. La comprensión de conceptos como la degeneración (donde un mismo nivel de energía corresponde a varios estados), la ortonormalización de bases y la interpretación probabilística de Born, constituye la base necesaria para cualquier desarrollo tecnológico o teórico en la física actual \citep{born1926}.

\section{Planteamiento del Trabajo}

Este Trabajo tiene como objetivo principal el análisis de un Hamiltoniano definido en un espacio de dimensión $d=4$. El problema consiste en determinar cómo una matriz hermítica puede ser descompuesta para revelar la física subyacente del sistema.

Se espera que este Trabajo contribuya a clarificar la relación entre los postulados teóricos abstractos y su aplicación práctica en la resolución de problemas dinámicos.


\section{Estructura del Trabajo}

Para lograr una exposición clara y autoconsistente, el Trabajo se ha estructurado de la siguiente manera:

El trabajo comienza con el Resumen para ofrecer una visión global de los hallazgos. El presente \textbf{Capítulo 1} constituye la Introducción, donde se contextualiza y justifica el Trabajo.

A continuación, el \textbf{Capítulo 2} detalla los \textbf{Objetivos} perseguidos. El \textbf{Capítulo 3} presenta el \textbf{Enunciado} del problema, definiendo el Hamiltoniano y el estado inicial objeto de estudio. Por su parte, el \textbf{Capítulo 4} recoge el \textbf{Desarrollo} técnico, donde se resuelven los autovalores, se construyen los proyectores, se calcula la evolución temporal y se analizan las constantes de movimiento.

Finalmente, se presenta la \textbf{Bibliografía} con las fuentes consultadas.


\chapter{Objetivos}
Este Trabajo busca desarollar la capacidad de trabajar con la descomposición espectral de un operador hermítico y aplicarla para analizar la evolución temporal de un estado cuántico en representación de Schrödinger. Se busca reforzar la comprensión formal de la relación entre hamiltonianos, operadores de evolución y observables medibles.

\chapter{Enunciado del problema}

El propósito de este Trabajo es el estudio cinemático y espectral de un sistema cuántico unidimensional definido en un espacio de Hilbert $\mathcal{H}$ de dimensión finita ($d=4$). El sistema está caracterizado por un operador Hamiltoniano $\hat{H}$ cuya representación matricial en una base ortonormal $\{|1\rangle, |2\rangle, |3\rangle, |4\rangle\}$ es:

\begin{equation}
\mathbf{H} = \begin{pmatrix} 
2 & 0 & 0 & 0 \\ 
0 & 1 & 1 & 0 \\ 
0 & 1 & 1 & 0 \\ 
0 & 0 & 0 & 3 
\end{pmatrix}
\end{equation}

A través del análisis de este sistema, se pretenden resolver los siguientes objetivos:

\begin{itemize}
    \item \textbf{Análisis Espectral:} Determinar la descomposición espectral del Hamiltoniano. Esto implica el cálculo de los autovalores de energía y la obtención de sus correspondientes proyectores ortogonales sobre los autoespacios del sistema.

    \item \textbf{Dinámica Cuántica:} Estudiar la evolución temporal de un estado inicial preparado específicamente como una superposición de estados de la base:

    \begin{equation}
    |\psi(0)\rangle = \frac{1}{\sqrt{2}} (|2\rangle + |3\rangle)
    \end{equation}
    Se debe obtener la expresión analítica del vector de estado $|\psi(t)\rangle$ para cualquier instante de tiempo $t > 0$ mediante la aplicación del operador de evolución unitaria $U(t) = e^{-iHt/\hbar}$.
    \item \textbf{Cálculo de Probabilidades:} Evaluar la probabilidad de transición o de medida, determinando la probabilidad de hallar al sistema en el autoestado $|4\rangle$ en un tiempo arbitrario $t$.

    \item \textbf{Constantes de Movimiento y Simetrías:} Investigar la existencia de observables compatibles $\hat{A}$ que conmuten con el Hamiltoniano ($[\hat{A}, \hat{H}] = 0$) y que posean al estado inicial $|\psi(0)\rangle$ como autovector.

\end{itemize}
\chapter{Desarrollo del Trabajo}

Se comienza abordando el análisis de un sistema cuántico unidimensional definido en un espacio de Hilbert $\mathcal{H}$ de dimensión $d=4$. El Hamiltoniano del sistema, representado por una matriz hermítica $H$, rige la dinámica y el espectro energético del sistema.

\section{Descomposición Espectral de $H$}

Un operador hermítico $H$ en un espacio de dimensión finita puede descomponerse en términos de sus autovalores $E_n$ y sus proyectores espectrales $P_n$ mediante la relación \citep{vonneumann1932}:
\begin{equation}
    H = \sum_{n=1}^{4} E_n P_n
\end{equation}


\subsection{Cálculo de Autovalores y Autovectores}

Para hallar los niveles de energía, se resuelve la ecuación característica $\det(H - \lambda I) = 0$ \citep{vonneumann1932}.

\begin{equation}
\begin{aligned}
\det(H - \lambda I) &= \begin{vmatrix} 2 - \lambda & 0 & 0 & 0 \\ 0 & 1 - \lambda & 1 & 0 \\ 0 & 1 & 1 - \lambda & 0 \\ 0 & 0 & 0 & 3 - \lambda \end{vmatrix} \\
  &= (2 - \lambda) (3 - \lambda) \left[(1 - \lambda)^2 - 1\right] \\
  &= (2 - \lambda) (3 - \lambda) \left[\lambda (\lambda-2) \right] = 0
\end{aligned}
\end{equation}

Tras el cálculo algebraico sobre la matriz dada, se obtienen los siguientes autovalores:

\begin{itemize}
    \item $E_1 = 2$
    \item $E_2 = 3$
    \item $E_3 = 0$
\end{itemize}

Dado que la solución $E_1 = 2$ es doble, es un sistema degenerado \citep{vonneumann1932}.

Posteriormente, Se determinan los autovectores correspondientes resolviendo el sistema $(H - E_n\mathbb{I})\ket{v_n} = 0$ \citep{dirac1930}. Tras el proceso de ortonormalización, la base de autovetores resultante es:

\begin{itemize}
  \item Para $E_1 = 2$:
  Al sustituir en el sistema $(H - 2\mathbb{I})\ket{v} = 0$, se obtiene la matriz:
  \begin{equation}
  \begin{pmatrix} 0 & 0 & 0 & 0 \\ 0 & -1 & 1 & 0 \\ 0 & 1 & -1 & 0 \\ 0 & 0 & 0 & 1 \end{pmatrix} \begin{pmatrix} x \\ y \\ z \\ w \end{pmatrix} = \begin{pmatrix} 0 \\ 0 \\ 0 \\ 0 \end{pmatrix}
  \end{equation}
  Esto permite definir dos autovectores ortonormales para este subespacio degenerado:
  \begin{equation}
  \ket{v_{2a}} = \begin{pmatrix} 1 \\ 0 \\ 0 \\ 0 \end{pmatrix}, \quad \ket{v_{2b}} = \frac{1}{\sqrt{2}} \begin{pmatrix} 0 \\ 1 \\ 1 \\ 0 \end{pmatrix}
  \end{equation}

  \item Para $E_2 = 3$:
  El sistema $(H - 3\mathbb{I})\ket{v} = 0$ resulta en:
  \begin{equation}
  \begin{pmatrix} -1 & 0 & 0 & 0 \\ 0 & -2 & 1 & 0 \\ 0 & 1 & -2 & 0 \\ 0 & 0 & 0 & 0 \end{pmatrix} \begin{pmatrix} x \\ y \\ z \\ w \end{pmatrix} = \begin{pmatrix} 0 \\ 0 \\ 0 \\ 0 \end{pmatrix}
  \end{equation}
  El autovector normalizado asociado es:
  \begin{equation}
  \ket{v_3} = \begin{pmatrix} 0 \\ 0 \\ 0 \\ 1 \end{pmatrix}
  \end{equation}

  \item Para $E_3 = 0$:
  El sistema $H\ket{v} = 0$ se representa mediante la matriz:
  \begin{equation}
  \begin{pmatrix} 2 & 0 & 0 & 0 \\ 0 & 1 & 1 & 0 \\ 0 & 1 & 1 & 0 \\ 0 & 0 & 0 & 3 \end{pmatrix} \begin{pmatrix} x \\ y \\ z \\ w \end{pmatrix} = \begin{pmatrix} 0 \\ 0 \\ 0 \\ 0 \end{pmatrix}
  \end{equation}
  Imponiendo la condición de normalización, el autovector es:
  \begin{equation}
  \ket{v_0} = \frac{1}{\sqrt{2}} \begin{pmatrix} 0 \\ 1 \\ -1 \\ 0 \end{pmatrix}
  \end{equation}
\end{itemize}
\subsection{Cálculo de los Proyectores Espectrales}

Una vez obtenida la base ortonormal de autovectores $\{\ket{v_{2a}}, \ket{v_{2b}}, \ket{v_3}, \ket{v_0}\}$, se procede a calcular los proyectores asociados a cada autoestado mediante el producto exterior $P_i = \ket{v_i}\bra{v_i}$. Estos operadores satisfacen las propiedades de idempotencia ($P_i^2 = P_i$) y hermiticidad ($P_i^\dagger = P_i$)\citep{dirac1930}.

Los proyectores asociados al autovalor degenerado $E_1 = 2$ son:
\begin{equation}
P_{2a} = \ket{v_{2a}}\bra{v_{2a}} = 
\begin{pmatrix} 1 \\ 0 \\ 0 \\ 0 \end{pmatrix}
\begin{pmatrix} 1 & 0 & 0 & 0 \end{pmatrix} =
\begin{pmatrix} 
1 & 0 & 0 & 0 \\ 
0 & 0 & 0 & 0 \\ 
0 & 0 & 0 & 0 \\ 
0 & 0 & 0 & 0 
\end{pmatrix}
\end{equation}

\begin{equation}
P_{2b} = \ket{v_{2b}}\bra{v_{2b}} = \frac{1}{2}
\begin{pmatrix} 0 \\ 1 \\ 1 \\ 0 \end{pmatrix}
\begin{pmatrix} 0 & 1 & 1 & 0 \end{pmatrix} =
\begin{pmatrix} 
0 & 0 & 0 & 0 \\ 
0 & 1/2 & 1/2 & 0 \\ 
0 & 1/2 & 1/2 & 0 \\ 
0 & 0 & 0 & 0 
\end{pmatrix}
\end{equation}

Para los autovalores simples $E_2 = 3$ y $E_3 = 0$, se obtiene:
\begin{equation}
P_{3} = \ket{v_{3}}\bra{v_{3}} = 
\begin{pmatrix} 0 \\ 0 \\ 0 \\ 1 \end{pmatrix}
\begin{pmatrix} 0 & 0 & 0 & 1 \end{pmatrix} =
\begin{pmatrix} 
0 & 0 & 0 & 0 \\ 
0 & 0 & 0 & 0 \\ 
0 & 0 & 0 & 0 \\ 
0 & 0 & 0 & 1 
\end{pmatrix}
\end{equation}

\begin{equation}
P_{0} = \ket{v_{0}}\bra{v_{0}} = \frac{1}{2}
\begin{pmatrix} 0 \\ 1 \\ -1 \\ 0 \end{pmatrix}
\begin{pmatrix} 0 & 1 & -1 & 0 \end{pmatrix} =
\begin{pmatrix} 
0 & 0 & 0 & 0 \\ 
0 & 1/2 & -1/2 & 0 \\ 
0 & -1/2 & 1/2 & 0 \\ 
0 & 0 & 0 & 0 
\end{pmatrix}
\end{equation}

\subsubsection{Verificación de los resultados}

Para validar la consistencia del desarrollo, se comprueba que se cumple la \textbf{Relación de Clausura}, la cual establece que la suma de los proyectores de una base completa debe ser igual a la identidad \citep{vonneumann1932}:
\begin{equation}
\sum_{i} P_i = P_{2a} + P_{2b} + P_3 + P_0 = 
\begin{pmatrix} 
1 & 0 & 0 & 0 \\ 
0 & 1 & 0 & 0 \\ 
0 & 0 & 1 & 0 \\ 
0 & 0 & 0 & 1 
\end{pmatrix} = \mathbb{I}
\end{equation}

Asimismo, se verifica la \textbf{Descomposición Espectral} del Hamiltoniano:
\begin{equation}
H = \sum_{i} E_i P_i = 2(P_{2a} + P_{2b}) + 3P_3 + 0P_0
\end{equation}
Al realizar la suma ponderada, se recupera la matriz original $H$, confirmando que los autovalores y autovectores calculados son correctos.

\section{Evolución Temporal del Sistema}

La dinámica del sistema viene descrita por la ecuación de Schrödinger. Para un Hamiltoniano independiente del tiempo, el estado del sistema en un instante $t$ se obtiene mediante la aplicación del operador de evolución unitaria $U(t) = e^{-iHt/\hbar}$ sobre el estado inicial $\ket{\psi(0)}$\citep{schrodinger1926}.

Se considera un estado inicial definido por una superposición equitativa de los autoestados asociados a las energías $E=2$ (tomando la dirección $\ket{v_{2b}}$) y $E=3$:
\begin{equation}
    \ket{\psi(0)} = \frac{1}{\sqrt{2}} (\ket{v_{2b}} + \ket{v_3})
\end{equation}

Dado que la acción del operador de evolución sobre un autoestado $\ket{v_n}$ se reduce a una fase multiplicativa $e^{-iE_n t/\hbar}$, el estado en un tiempo arbitrario $t$ resulta:
\begin{equation}
    \ket{\psi(t)} = \frac{1}{\sqrt{2}} \left( e^{-i 2t/\hbar} \ket{v_{2b}} + e^{-i 3t/\hbar} \ket{v_3} \right)
\end{equation}

Sustituyendo los vectores columna calculados anteriormente, el vector de estado evoluciona como:
\begin{equation}
    \ket{\psi(t)} = \frac{1}{\sqrt{2}} \begin{pmatrix} 0 \\ \frac{1}{\sqrt{2}}e^{-i 2t/\hbar} \\ \frac{1}{\sqrt{2}}e^{-i 2t/\hbar} \\ e^{-i 3t/\hbar} \end{pmatrix}
\end{equation}

\section{Cálculo de Probabilidades}

Un postulado fundamental de la mecánica cuántica establece que la probabilidad de encontrar el sistema en un estado $\ket{\phi}$ tras realizar una medida en el instante $t$ viene dada por la regla de Born \citep{born1926}:
\begin{equation}
    \mathcal{P}_{\phi}(t) = |\braket{\phi}{\psi(t)}|^2
\end{equation}

En este caso, se requiere determinar la probabilidad de hallar el sistema en el estado $\ket{4}$, que corresponde al cuarto vector de la base canónica del espacio de Hilbert, $\ket{4} = (0, 0, 0, 1)^T$. Se observa que, según el cálculo previo, este vector coincide exactamente con el autovector $\ket{v_3}$.

La amplitud de probabilidad es:
\begin{equation}
    \braket{4}{\psi(t)} = \frac{1}{\sqrt{2}} \left( e^{-i 2t/\hbar} \braket{4}{v_{2b}} + e^{-i 3t/\hbar} \braket{4}{v_3} \right)
\end{equation}

Debido a la ortogonalidad de la base de autovectores, tenemos que $\braket{4}{v_{2b}} = 0$ y $\braket{4}{v_3} = 1$. Por lo tanto:
\begin{equation}
    \braket{4}{\psi(t)} = \frac{1}{\sqrt{2}} e^{-i 3t/\hbar}
\end{equation}

Finalmente, la probabilidad resulta ser:
\begin{equation}
    \mathcal{P}_4(t) = \left| \frac{1}{\sqrt{2}} e^{-i 3t/\hbar} \right|^2 = \frac{1}{2} |e^{-i 3t/\hbar}|^2 = \frac{1}{2}
\end{equation}

Este resultado indica que la probabilidad de encontrar el sistema en el estado $\ket{4}$ es del $50\%$ y permanece constante en el tiempo, ya que el estado inicial está compuesto por autoestados estacionarios y la medida coincide con uno de los elementos de dicha base.
\section{Constantes de Movimiento y Simetrías}

Un observable $\hat{A}$ se define como una constante de movimiento si conmuta con el Hamiltoniano del sistema, es decir, $[\hat{A}, \hat{H}] = \hat{A}\hat{H} - \hat{H}\hat{A} = 0$. Esta condición implica que ambos operadores comparten una base común de autovectores y que el valor esperado de $\hat{A}$ permanece constante en el tiempo \citep{noether1918}.

\subsection{Búsqueda de un Observable Compatible}

Se requiere investigar un observable $\hat{A}$ que, además de conmutar con $\hat{H}$, posea al estado inicial $\ket{\psi(0)} = \frac{1}{\sqrt{2}}(\ket{v_{2b}} + \ket{v_3})$ como autovector. 

Para que $\ket{\psi(0)}$ sea autovector de $\hat{A}$ con autovalor $a$, se debe cumplir:
\begin{equation}
    \hat{A} \ket{\psi(0)} = a \ket{\psi(0)}
\end{equation}

Dado que $[\hat{A}, \hat{H}] = 0$, ambos operadores deben compartir una base común de autovectores, lo que permite la determinación simultánea de sus valores propios \citep{heisenberg1927}. Un candidato natural para este observable es el propio proyector sobre el estado inicial, $\hat{A} = \hat{P}_{\psi(0)}$. Sin embargo, para que este proyector sea una constante de movimiento, debe ser funcionalmente dependiente del Hamiltoniano o proyectar sobre un autoespacio degenerado \citep{vonneumann1932}.

Otra opción relevante es considerar el operador paridad o un operador de simetría interna si el Hamiltoniano presenta bloques idénticos. En este caso, dado que $\ket{\psi(0)}$ es una combinación de estados con energías $E=2$ y $E=3$, un observable compatible sería aquel que asigne un valor constante a dicho subespacio, como por ejemplo:
\begin{equation}
    \hat{A} = 2(P_{2a} + P_{2b}) + 3P_3
\end{equation}

% RELLENAR: Inserta aquí los vectores columna

\begin{thebibliography}{99}

\bibitem[Planck(1900)]{planck1900} 
Planck, M. (1900). Zur Theorie des Gesetzes der Energieverteilung im Normalspectrum. \emph{Verhandlungen der Deutschen Physikalischen Gesellschaft}, 2, 237-245.

\bibitem[Dirac(1930)]{dirac1930} 
Dirac, P. A. M. (1930). \emph{The Principles of Quantum Mechanics}. Oxford University Press.

\bibitem[Schrödinger(1926)]{schrodinger1926} 
Schrödinger, E. (1926). Quantisierung als Eigenwertproblem. \emph{Annalen der Physik}, 384(4), 361-376.

\bibitem[Feynman(1982)]{feynman1982} 
Feynman, R. P. (1982). Simulating physics with computers. \emph{International Journal of Theoretical Physics}, 21(6), 467-488.

\bibitem[Born(1926)]{born1926} 
Born, M. (1926). Zur Quantenmechanik der Stoßvorgänge. \emph{Zeitschrift für Physik}, 37(12), 863-867.

\bibitem[Von Neumann(1932)]{vonneumann1932} 
Von Neumann, J. (1932). \emph{Mathematische Grundlagen der Quantenmechanik}. Springer-Verlag.

\bibitem[Noether(1918)]{noether1918} 
Noether, E. (1918). Invariante Variationsprobleme. \emph{Nachrichten von der Gesellschaft der Wissenschaften zu Göttingen}, 1918, 235-257.

\bibitem[Heisenberg(1927)]{heisenberg1927} 
Heisenberg, W. (1927). Über den anschaulichen Inhalt der quantentheoretischen Kinematik und Mechanik. \emph{Zeitschrift für Physik}, 43(3-4), 172-198.
\end{thebibliography}
%\bibliographystyle{plain} 
%\bibliography{bibliografia}

\end{document}

